
\chapter{结论}

\section{总结}

本文围绕模板配置、正文写作和常用可选功能，介绍了如何利用 \LaTeX 撰写南开大学学位论文。通过以上示例，可以总结出以下关键要点：

\begin{enumerate}
    \item 选择合适的模板是论文撰写的基础，南开大学 \LaTeX 模板提供了规范的格式支持
    \item 合理组织文档结构，使用模块化的文件管理方式可以提高写作效率
    \item 掌握基本的 \LaTeX 语法和学术写作功能是高质量论文的保障
    \item 合理使用打印模式、匿名评审页、附录等可选功能，可以减少后期格式调整成本
    \item 注意字体兼容性和编译环境配置，确保文档在不同平台上的一致性
\end{enumerate}

\section{使用建议}

对于准备使用 \LaTeX 撰写南开大学学位论文的同学，建议：

\begin{itemize}
    \item 提前熟悉 \LaTeX 基本语法，可以通过在线教程或示例文档学习
    \item 使用可靠的编译平台，如 TeXPage 或 Overleaf，避免环境配置问题
    \item 参考模板提供的示例，逐步掌握各种学术元素的排版方法
    \item 定期备份文档，避免因编译错误导致的工作丢失
\end{itemize}

通过合理使用 \LaTeX 和南开大学学位论文模板，用户可以将更多精力集中在研究内容本身，而不是重复性的格式调整工作。
